\documentclass[11pt]{article}

%======================= Packages =======================
\usepackage[margin=1in]{geometry}
\usepackage{amsmath,amssymb,amsthm,mathtools}
\usepackage{bm}
\usepackage{enumitem}
\usepackage{graphicx}
\usepackage{float}
\usepackage{booktabs}
\usepackage{xcolor}
\usepackage{hyperref}
\hypersetup{colorlinks=true,linkcolor=blue,citecolor=blue,urlcolor=blue}

\graphicspath{{../figures/}{figures/}}

%======================= Math macros =======================
\newcommand{\E}{\mathbb{E}}
\newcommand{\Var}{\mathrm{Var}}
\newcommand{\Cov}{\mathrm{Cov}}
\newcommand{\Fset}{\mathcal{F}}
\newcommand{\R}{\mathbb{R}}
\newcommand{\Id}{\mathbb{1}}
\newcommand{\tp}{^{\prime}}
\newcommand{\eps}{\varepsilon}
\newcommand{\beps}{\bm{\varepsilon}}
\newcommand{\Mmap}{\mathcal{M}}
\newcommand{\Om}{\Omega}
\newcommand{\pinv}[1]{#1^{+}}
\DeclareMathOperator{\MSPE}{MSPE}
\DeclareMathOperator{\RMSE}{RMSE}
\DeclareMathOperator{\MAE}{MAE}
\DeclareMathOperator{\CRPS}{CRPS}
\DeclareMathOperator{\chol}{chol}
\DeclareMathOperator{\rank}{rank}
\DeclareMathOperator{\diag}{diag}
\DeclareMathOperator{\vecop}{vec}

\theoremstyle{remark}
\newtheorem{remark}{Remark}
\newtheorem*{note}{Note}

%---- figure helper: include the real PNG (resolved via \graphicspath: figures/ or ../figures/)
\newcommand{\projfig}[3]{% {relpath}{caption}{label}
  \begin{figure}[H]\centering
  \includegraphics[width=0.86\linewidth]{#1}%
  \caption{#2}\label{#3}
  \end{figure}}

\title{Conditional Forecasting with Uncertainty:\\
Scenarios, Soft Conditions, and Historical Counterfactuals\\
in a Monetary VAR\\[4pt]
\large A methods-and-tooling companion}
\author{%
  Etienne Latulippe\\
  \small Built on the Canadian monetary VAR of Moran, Stevanovic \& Surprenant (BoC SWP 2025-28)}
\date{\today}

\begin{document}
\maketitle

\begin{abstract}
\noindent Central-bank risk-scenario exercises are typically reported as point
paths: ``if oil rises to \$120, inflation does this.'' The underlying
Waggoner--Zha (WZ) conditional-forecast machinery, however, is fundamentally a
statement about a \emph{distribution} of shocks consistent with a scenario, and
it therefore carries its own measure of uncertainty that is almost never shown.
This paper presents an operational toolkit that puts that uncertainty front and
centre. Working in a reduced-form vector autoregression identified recursively,
we (i) derive the conditional-forecast \emph{fan chart}, decomposing its width
into a closed-form \emph{shock} (innovation) component and a simulated
\emph{parameter} (estimation) component, and show that the constrained variable's
band collapses to zero by construction; (ii) formalize \emph{soft} conditions,
which condition on a distribution for the constrained variable rather than an
exact path, so that the constrained variable itself acquires a band; (iii) treat
\emph{joint} and \emph{partial-horizon} scenarios as nothing more than extra rows
of one linear restriction; and (iv) derive \emph{historical counterfactuals} as a
minimum-norm perturbation of realized history, optionally confined to a chosen
structural shock. Each object is stated with the intuition and the exact algebra
implemented in the accompanying code, with full derivations in the appendix. The
toolkit is delivered as a reproducible offline pipeline and a self-contained
interactive browser dashboard in which the impulse responses, scenario fan
charts, and counterfactuals are recomputed live. The Canadian monetary VAR of
Moran, Stevanovic \& Surprenant (2025) serves as the running validation case.
\end{abstract}

%=========================================================
\section{Introduction}
%=========================================================
Scenario analysis is a staple of policy institutions: condition the model on a
hypothetical path for one or more variables and read off what the rest of the
economy does. The dominant technical device is the Waggoner--Zha (1999)
conditional forecast, which finds the ``most likely'' sequence of future shocks
consistent with the imposed path and propagates it to every variable. Moran,
Stevanovic \& Surprenant (2025, hereafter MSS) apply exactly this device to a
monthly vector autoregression (VAR) for the Canadian economy and report four risk
scenarios. Those exercises, like most of their kind, report conditional
\emph{means}. Yet Waggoner and Zha's original contribution is largely about the
\emph{distribution} around those means. A scenario is produced by shocks, and the
shocks it does not pin down remain random, so the conditional forecast of every
unconstrained variable is itself a random object with a well-defined predictive
distribution. Reporting only its mean discards precisely the information a risk
exercise is meant to convey---how uncertain the consequences of the scenario are.

This paper is a methods-and-tooling companion that makes conditional forecasting
operational \emph{with} its uncertainty. Its central object is the
conditional-forecast fan chart, whose width we derive and decompose into two
additive sources: a \emph{shock} component---the innovations the constraint
leaves free, available in closed form---and a \emph{parameter} component---the
estimation uncertainty in the VAR, obtained by re-solving the scenario over
posterior or bootstrap draws of the coefficients. A by-product of the derivation,
visible in every figure below, is that the constrained variable's band collapses
to a point by construction: a scenario pins what it pins with certainty and lets
uncertainty accumulate only where the model is free to move. Around this core we
develop the vocabulary a working risk desk actually needs. Soft conditions
replace an exact path by a distribution---``oil around \$120, give or take''---so
that the constrained variable, too, carries a band; joint and partial-horizon
scenarios, which constrain several variables at once or a variable over only part
of the horizon, turn out to be nothing more than adding or deleting rows of a
single linear restriction; and historical counterfactuals---``what would
inflation have done had the Bank tightened 100 basis points more?''---emerge as a
minimum-norm perturbation of realized history, optionally confined to a single
structural shock. Each object is stated with its intuition and the exact algebra
implemented in the accompanying code, and derived in full in the appendices.

The apparatus is delivered as software, and the delivery is part of the point.
The estimation runs offline from a cached panel, so a result is deterministic and
archivable, and it is guarded by a regression suite that checks the algebraic
invariants any correct implementation must satisfy. A single self-contained
browser dashboard then exposes the toolkit interactively: the impulse responses,
the scenario fan charts with their live shock band, and the counterfactuals are
all recomputed in the browser from the exported model objects, with no server and
no re-run of the estimation. Throughout we take the estimated VAR and its
recursive identification from MSS as given and validated; the reduced-form model,
the baseline forecast, the conditional-forecast mean, and the variance
decomposition are reproduced, for completeness, in Appendix~\ref{app:orig}, and
the new derivations---the uncertainty decomposition, soft and joint conditions,
and counterfactuals---are collected in Appendix~\ref{app:new}. The Canadian
monetary VAR of MSS is the running example, and the figures are those produced by
the pipeline and archived in the HTML report that accompanies the project.

Notation is kept light. $y_t\in\R^{K}$ is the vector of endogenous variables, $p$
the lag order, $H$ the horizon, and $\Fset_t=\{y_s\}_{s\le t}$ the information set.
Reduced-form innovations $u_t$ have mean zero and covariance $\Sigma_u$; a
structural impact matrix $D$ with $\Sigma_u=DD\tp$ maps them to orthonormal
structural shocks through $u_t=D\eps_t$, $\E[\eps_t\eps_t\tp]=I_K$. Reduced-form
moving-average coefficients are $\Phi_j=JA^jJ\tp$ (with $A$ the companion matrix
and $J$ the top-block selector), structural ones $\Theta_j=\Phi_jD$, and
$\pinv{R}=R\tp(RR\tp)^{-1}$ is the pseudoinverse of a full-row-rank $R$.

%=========================================================
\section{The model and conditional forecasts}\label{sec:model}
%=========================================================
We work with the reduced-form VAR
\begin{equation}
y_t=\nu+A_1y_{t-1}+\dots+A_py_{t-p}+u_t,
\label{eq:var}
\end{equation}
written in companion form $Y_t=\mu+AY_{t-1}+U_t$ and iterated forward. Taking the
conditional expectation kills every future shock, leaving the baseline
(unconditional) forecast
\begin{equation}
\E[y_{t+h}\mid\Fset_t]=\Big(\textstyle\sum_{j=0}^{h-1}JA^{j}J\tp\Big)\nu+JA^{h}Y_t,
\label{eq:baseline}
\end{equation}
while the deviation of any realized path from this baseline is driven entirely by
future shocks,
\begin{equation}
y_{t+h}-\E[y_{t+h}\mid\Fset_t]=\sum_{j=0}^{h-1}\Phi_j\,u_{t+h-j}
=\sum_{j=0}^{h-1}\Theta_j\,\eps_{t+h-j}.
\label{eq:dev}
\end{equation}
These steps are standard and Appendix~\ref{app:orig} records them in full; the one
object the rest of the paper needs is the stacked structural map. Collecting
\eqref{eq:dev} for $h=1,\dots,H$ and writing
$\beps=[\eps_{t+1}\tp,\dots,\eps_{t+H}\tp]\tp$ for the stack of future structural
shocks, the deviation of the whole forecast path from the baseline is a single
linear image of those shocks,
\begin{equation}
\Delta y=\Mmap\,\beps,
\qquad
\Mmap=\begin{bmatrix}
\Theta_0 & & & \\
\Theta_1 & \Theta_0 & & \\
\vdots & & \ddots & \\
\Theta_{H-1} & \Theta_{H-2} & \cdots & \Theta_0
\end{bmatrix},
\label{eq:Mmap}
\end{equation}
a block-lower-triangular $KH\times KH$ matrix with $\Theta_0=D$ on the diagonal.
Everything below is a linear operation on \eqref{eq:Mmap}; working in the
structural stack, where $\beps$ has identity covariance, is what keeps the
uncertainty algebra of the next section clean.

A scenario imposes target values on a subset of variables over some horizons.
Each imposed cell is one linear equation in $\beps$, obtained by selecting the
corresponding row of $\Mmap$, so a scenario is the linear restriction
\begin{equation}
R\,\beps=r,\qquad R=\Mmap_{[\text{constrained rows}],\,:},\quad
r=\big(\text{target}-\text{baseline}\big)_{[\text{constrained rows}]},
\label{eq:restriction}
\end{equation}
with $r\in\R^{m}$ collecting the $m$ constrained cells. Because a scenario fixes
far fewer cells than there are free shocks, \eqref{eq:restriction} is
underdetermined; Waggoner and Zha, after Doan, Litterman and Sims, select the
minimum-norm solution
\begin{equation}
\hat\beps=\arg\min_{\beps}\tfrac12\beps\tp\beps \ \ \text{s.t.}\ \ R\beps=r
=R\tp(RR\tp)^{-1}r=\pinv{R}\,r,
\label{eq:wzsol}
\end{equation}
and the conditional forecast is $\Delta y=\Mmap\hat\beps$ added back to the
baseline. Two points recur below. First, minimizing $\beps\tp\beps$ minimizes the
sum of squared \emph{structural} shocks, which is the Mahalanobis norm
$u\tp\Sigma_u^{-1}u$ of the reduced-form innovations and, under $\beps\sim
N(0,I)$, picks out the mode of the shocks consistent with the scenario; minimizing
the unweighted reduced-form norm instead would load implausibly large shocks, in
standard-deviation terms, onto low-variance equations and distort every
unconstrained path, so the distinction is substantive rather than cosmetic.
Second, since $\E\lVert\beps\rVert^2=KH$, the ratio $\lVert\hat\beps\rVert/
\sqrt{KH}$ is a readable gauge of how demanding a scenario is---near one for an
ordinary draw, several times that for a scenario the model regards as close to
impossible---and the dashboard reports it as a traffic light on every exercise.

%=========================================================
\section{Scenarios and their uncertainty}\label{sec:scen}\label{sec:unc}
%=========================================================
The restriction \eqref{eq:restriction} is agnostic about how a target path is
chosen. In practice a scenario is described in economic terms and translated into
the VAR's transformed units: a variable in levels (the unemployment or policy
rate) is constrained directly, while a variable in log-differences (oil, prices,
housing, employment, the exchange rate) has its authored \emph{level} path
$\ell_{1:H}$ mapped by
\begin{equation}
\text{level variable: } r^{\text{unit}}_h=\ell_h,
\qquad
\text{log-diff variable: } r^{\text{unit}}_h=\log\ell_h-\log\ell_{h-1},
\label{eq:units}
\end{equation}
with $\ell_0$ the last observed level. A small vocabulary of shapes---a
\emph{ramp} to a target and hold, a flat \emph{hold}, a step--hold--ease
\emph{policy} rate path, a \emph{replay} of a variable's own historical growth
from a chosen date, and a \emph{pulse} that ramps, holds, and is then
released---covers the common cases; the constructors are spelled out in
Appendix~\ref{app:scenmap}.

A scenario, however, is not a single line. The constraint fixes only $m$ linear
combinations of the $KH$-dimensional shock vector and leaves the rest free, so the
conditional forecast of every unconstrained variable has a genuine predictive
distribution. Holding the parameters fixed, conditioning the prior $\beps\sim
N(0,I)$ on $R\beps=r$ gives another Gaussian, supported on the constraint
subspace,
\begin{equation}
\beps\mid R\beps=r\ \sim\ N(\hat\beps,\,P),\qquad
\hat\beps=\pinv{R}r,\quad P=I-\pinv{R}R,
\label{eq:shockpost}
\end{equation}
where $P$ is the orthogonal projector onto the null space of $R$; propagating
through \eqref{eq:Mmap}, the conditional forecast is
$\Delta y\sim N(\Mmap\hat\beps,\ \Mmap P\Mmap\tp)$,
\begin{equation}
\Var(\Delta y\mid R\beps=r)=\Mmap P\Mmap\tp.
\label{eq:shockband}
\end{equation}
Because $RP=0$, the rows of $\Mmap P\Mmap\tp$ that belong to constrained cells
vanish: the fan chart pinches to a point on the variable and horizons the scenario
pins, and opens up only where the model is free. This \emph{shock} band is closed
form---no simulation is needed---which is what lets the dashboard draw it live as
the user edits a scenario. Since the display transforms are nonlinear (levels are
reconstructed as $\ell_0\exp(\text{cumsum})$, growth rates as year-over-year
differences), bands are read as sample quantiles of transformed draws from
\eqref{eq:shockpost}, not as transforms of Gaussian quantiles;
Appendix~\ref{app:shock} gives the argument.

The estimated parameters carry their own uncertainty, and it cannot be bolted on
afterward: the coefficients enter $R$, $r$, and $\Mmap$ themselves, so a different
draw of the VAR implies a different constraint system and a different conditional
mean, and the scenario must be re-solved inside each draw. Drawing from the
normal--inverse--Wishart posterior when a shrinkage prior is used, or from a
bias-adjusted residual bootstrap otherwise, and taking quantiles of the
transformed conditional draws across parameter draws yields the \emph{parameter}
band; combined with the shock band it gives the reported fan
(Appendix~\ref{app:param}). One reading caveat is worth stating in the main text:
the point baseline from the OLS estimate is not the centre of the bootstrap
distribution, because the bias adjustment and the rejection of explosive draws
shift the per-draw baseline, so the clean object for ``is this scenario
distinguishable from the baseline?'' is the scenario-minus-baseline band computed
draw by draw, which nets out the shift.

Partial-horizon conditions come for free: leaving a horizon out of
\eqref{eq:restriction} leaves it free, and the pulse type does exactly this,
constraining a plateau and releasing everything after it. For a log-difference
variable, pinning months $1,\dots,n$ pins the \emph{level} at month $n$ and lets
the model supply the persistence, so ``oil spikes for six months and then
normalizes'' becomes a statement the model completes rather than a plateau the
analyst imposes. Figure~\ref{fig:temp} is such a released oil pulse. Over the six
constrained months the effect is imposed, but the tail is the model's own answer:
headline inflation climbs from about 2\% to a peak near 3.8\% in 2027 and then
eases back toward 2.5\%, housing inflation recovers from mildly negative toward
zero, and unemployment and the policy rate barely stir. The bands widen steadily,
and by 2028 the 90\% interval for inflation runs from roughly zero to 5\%---a
reminder that a six-month spike says little about the medium term with any
confidence, which is exactly what a risk exercise should make visible.

\projfig{temporary_shocks/figure5_bands_oil_spike_6_months.png}
{A temporary (``pulse'') oil scenario: oil is constrained for six months and then
released, so the tail is the model's own response. Green is history, red the
scenario median, and the shaded bands the 68\% and 90\% intervals combining shock
and parameter uncertainty.}
{fig:temp}

Constraining several variables at once merely adds rows to $R$ and $r$, so
genuinely multivariate exercises are as cheap as univariate ones. This is what
makes qualitatively new scenarios reachable: constraining oil up and unemployment
up together forces the system onto a supply-shock path that neither constraint
produces alone; holding activity down while the policy rate is fixed expresses a
``recession with no policy response'' as a joint condition; and a tariff shock,
which hits both foreign activity and the exchange rate, must constrain the two
together, since constraining only one lets the model offset it through the other.
Figure~\ref{fig:joint} shows the stagflation case. Unemployment is held at 8\%,
visible as the flat median with essentially no band---the zero-width pinch of
\eqref{eq:shockband} made concrete---while oil is pushed up; inflation rises
briefly toward 3\% and then drifts below 2\% as the demand drag dominates, and the
policy rate falls through the window with a wide band. The only quantity that
moves against the analyst is feasibility: a joint condition is far more demanding
than either leg, and $\lVert\hat\beps\rVert$ grows accordingly.

\projfig{joint_scenarios/figure5_bands_stagflation.png}
{A joint scenario (stagflation): oil and unemployment are constrained upward
together. The unemployment panel shows the constrained path as a flat median with
no band---the constrained variable's uncertainty is zero by construction---while
the free variables fan out.}
{fig:joint}

A hard condition pins the constrained variable to an exact number, but ``oil
around \$120'' is more honestly a distribution than a point. Writing the
restriction as a noisy measurement $r=R\beps+e$ with $e\sim N(0,\Om)$,
\begin{equation}
r=R\beps+e,\qquad \Om=\diag(\bm\sigma^2),
\label{eq:soft}
\end{equation}
and updating the prior in the usual Gaussian way gives the soft posterior
\begin{equation}
\beps\mid r\ \sim\ N\!\big(R\tp(RR\tp+\Om)^{-1}r,\ \ I-R\tp(RR\tp+\Om)^{-1}R\big),
\label{eq:softpost}
\end{equation}
which returns the hard solution as $\Om\to0$ and otherwise gives the constrained
variable a nondegenerate band of its own (Appendix~\ref{app:soft}). Softening the
same oil pulse, Figure~\ref{fig:soft}, barely moves the medians of the reported
variables and only modestly widens their near-term fans; the substantive change is
on the constrained variable itself, whose pinned line becomes a band---the more
truthful object for something labelled a risk scenario. Softness is set per
horizon, so a path can be pinned tightly early and loosely later.

\projfig{soft_conditions/figure5_bands_oil_spike_loosely_held.png}
{A soft oil condition: conditioning on a distribution rather than an exact path.
On the reported variables it is close to the hard pulse of
Figure~\ref{fig:temp}; the difference is that the constrained variable itself now
carries uncertainty.}
{fig:soft}

%=========================================================
\section{Historical counterfactuals}\label{sec:cf}
%=========================================================
Nothing in the algebra confines it to the forecast horizon. Run the same
machinery over a window the data actually realized and it answers historical
``what if'' questions. Over such a window from an origin $t_0$, write the outcome
as the origin baseline plus the realized innovations,
$y_{\text{actual}}=F_{\text{base}}+\Mmap\beps_{\text{actual}}$, so that on the
constrained rows $R\beps_{\text{actual}}=r_{\text{actual}}$ is the realized
deviation from the baseline. A counterfactual imposes a different constrained path
$r_{\text{cf}}$ while changing as little else as possible,
\begin{equation}
\min_{\delta}\ \lVert\delta\rVert\quad\text{s.t.}\quad
R(\beps_{\text{actual}}+\delta)=r_{\text{cf}}
\ \Longrightarrow\
\delta=\pinv{R}\,(r_{\text{cf}}-r_{\text{actual}}),
\label{eq:cf}
\end{equation}
so that
\begin{equation}
y_{\text{cf}}=y_{\text{actual}}+\Mmap\,\pinv{R}\,(r_{\text{cf}}-r_{\text{actual}}).
\label{eq:cf2}
\end{equation}
The baseline cancels: the counterfactual is a perturbation of \emph{history}, not
of a forecast, so it keeps every realized shock the constraint does not require
changing. A pure policy experiment should move history only through the
monetary-policy shock, which is imposed by restricting $\delta$ to the columns of
$R$ carrying that shock before solving \eqref{eq:cf} (Appendix~\ref{app:cf}).
Figure~\ref{fig:cf} runs precisely this exercise: had the Bank held the policy
rate 100 basis points higher through 2022--23, the model puts the price level
about 0.6 percentage points lower and the unemployment rate about 0.75 points
higher by the end of the window---the shaded wedge between the actual and
counterfactual paths---a disinflation bought with labour-market slack, and the
expected sign. Because the change is confined to the policy shock, the 2022
commodity and supply shocks remain in both paths and only the monetary leg
differs. Two caveats travel with any such number. The reduced-form VAR is assumed
invariant to the policy change contemplated, which is exactly what the Lucas
critique denies, so the exercise is informative about the model's internal
propagation, not about what the economy would truly have done; and the size of the
required perturbation, $\lVert\delta\rVert/\sqrt{KH}$, gauges how far the
counterfactual reaches beyond historical experience. Both are surfaced in the
dashboard.

\projfig{counterfactual/tighter_boc_price_unemp.png}
{A historical counterfactual: the price level and unemployment over 2022--23 under
a 100bp-tighter policy path (dashed) against the realized data (solid). The
end-of-window gaps are annotated; the perturbation is confined to the policy
shock.}
{fig:cf}

%=========================================================
\section{Identification}\label{sec:ident}
%=========================================================
Everything above is conditional on an identification of the impact matrix $D$. We
use the recursive scheme of the source paper, $D=\chol(\Sigma_u)$ lower
triangular, which is a set of \emph{timing restrictions}: because $D$ is lower
triangular, a shock ordered $k$ moves variables $1,\dots,k-1$ only with a lag and
variables $k,\dots,K$ within the period. The ordering therefore encodes a
contemporaneous causal chain, and the one used here---global and commodity
variables first, then the domestic price block, then activity and the policy
rate, with forward-looking financial variables last---reads as a set of economic
priors: block exogeneity of a small open economy to world conditions within the
month, prices that are sticky relative to the activity and policy block, and asset
prices that respond to everything contemporaneously. These are assumptions of
convenience as much as of theory; they are not tested by the data, and every
object that is expressed in terms of \emph{named} structural shocks inherits them.

It is worth being precise about which results depend on this choice, because it is
less than one might expect. The Waggoner--Zha conditional forecast minimises the
Mahalanobis norm $u\tp\Sigma_u^{-1}u$ of the reduced-form innovations, and both
the reduced-form dynamics and $\Sigma_u$ are invariant to how one factors
$\Sigma_u=DD\tp$. The scenario means of Section~\ref{sec:scen}, and the shock
bands built from the conditional covariance, are therefore \emph{invariant} to the
Cholesky ordering: permuting the recursive order, or indeed replacing the Cholesky
factor by any other valid $D$, leaves every conditional forecast and every fan
chart unchanged. What does depend on the identification is anything attributed to
a particular structural shock: the impulse responses $\Theta_j=\Phi_jD$, the
forecast-error variance decomposition, and---among the exercises developed
here---the counterfactual that is \emph{confined to a chosen shock}, such as the
tighter-policy experiment of Section~\ref{sec:cf}, whose confinement singles out
the policy column of $D$. Identification thus matters exactly where a result is
read as ``the effect of shock $\ell$,'' and is immaterial where a scenario is
expressed purely as a path for observed variables.

That distinction also bounds the cost of the recursive scheme's known weaknesses.
Recursive timing restrictions are prone to the familiar puzzles---a
contractionary monetary shock that appears to raise prices, or an activity
response of the wrong sign---and the results can be sensitive to the ordering,
which is hard to defend at the monthly frequency for variables that in truth
respond to one another without delay. Because the identification is a single
input, $D$, the toolkit is agnostic to how it is obtained, and a natural line of
further work is to replace the Cholesky factor with an identification that rests
on economics rather than timing. Several are under review: sign restrictions on
the impulse responses, which set-identify the shocks from theory-consistent
comovements (Uhlig, 2005; Rubio-Ram\'irez, Waggoner and Zha, 2010); external-
instrument or proxy-SVAR identification, which uses an outside series correlated
with a single structural shock---high-frequency policy surprises for the monetary
shock, for instance (Mertens and Ravn, 2013; Gertler and Karadi, 2015; Stock and
Watson, 2018); identification through heteroskedasticity, which exploits changes
in the shock variances across regimes (Rigobon, 2003); and narrative restrictions
that discipline the shocks with information about specific historical episodes
(Antol\'in-D\'iaz and Rubio-Ram\'irez, 2018). Each of these delivers a different
$D$ (or a set of admissible $D$'s), and each can be substituted into the pipeline
without touching the scenario machinery: the impulse responses, the variance
decomposition, and the shock-confined counterfactuals would update, while the
scenario forecasts and their bands would not move at all.

%=========================================================
\section{Implementation and the dashboard}\label{sec:tooling}
%=========================================================
The methods are delivered as software, and the software is part of the
contribution. The estimation runs from a cached panel with no network access, so a
run is deterministic and archivable, and a one-time online step freezes the data
into that cache. A regression suite checks the data-independent invariants any
correct refactor must preserve---that $\Sigma_u$ is symmetric positive definite,
that $DD\tp$ recovers it, that the impact response equals $D$, that
forecast-error-variance shares sum to one, that hard conditioning hits its
constraint to numerical zero, and that a counterfactual with
$r_{\text{cf}}=r_{\text{actual}}$ reproduces history exactly---alongside a
golden-value snapshot that catches numerical drift. A report generator collects
every figure and the variance-decomposition table into a single self-contained
HTML file.

The interactive dashboard is a single self-contained HTML page that does all of
its computation in the browser from the exported model objects: the structural
responses $\Theta_j$, the impact matrix $D$, the baseline, the companion matrix,
and the transformed history. An impulse-response explorer rescales and recumulates
responses on the fly; a scenario builder lets the user assemble hard or soft,
single or joint, full or partial-horizon constraints and watch the conditional
forecast---and, on demand, the closed-form shock band of
Section~\ref{sec:scen}---recompute as the inputs change, with the in-browser
Waggoner--Zha solve validated against the reference implementation to about
$10^{-8}$; a counterfactual tab implements \eqref{eq:cf2} with the
structural-shock confinement and the feasibility and Lucas-critique readouts; and
a grouped gallery holds the archived fan charts. Because the browser recomputes
from exported matrices rather than re-estimating, the parameter-uncertainty band
remains an offline object; exporting a set of posterior draws to make the full
band live is a natural next step.

%=========================================================
\section{Evaluation and outlook}\label{sec:oos}
%=========================================================
The apparatus is internally coherent and reproducible, and the figures behave as
the theory says they should---the constrained variable pinches, the free variables
fan out, soft conditions widen the constrained variable's own band, and the
counterfactual signs are the expected ones. What remains is to show that the
conditional forecasts are also \emph{good}. The natural test is a pseudo-real-time
evaluation: re-estimate on an expanding window and score the forecasts against
realized data, both by point-forecast loss relative to an autoregressive benchmark
and by a proper density score computed from predictive quantiles, which would also
tell us whether the fan-chart widths are calibrated. The evaluation framework---the
benchmark, the error definitions, and the quantile representation of the continuous
ranked probability score---is standard and is recorded in Appendix~\ref{app:orig};
the backtesting harness that would populate it is the principal item of future
work, and it reuses the same conditional-forecast solve developed here inside a
date loop.

To conclude, conditional forecasting is usually presented as a point exercise, but
the method is a statement about a distribution of shocks and therefore carries its
own uncertainty. Making that uncertainty the object of interest---a fan chart with
a closed-form shock component and a simulated parameter component, soft conditions
that give the constrained variable a band, joint and partial-horizon scenarios as
row operations on one restriction, and historical counterfactuals as minimum-norm
perturbations of history---turns a familiar tool into a more honest one, and the
interactive delivery puts it in the hands of anyone with a browser. The remaining
gap is empirical validation of the forecasts themselves, which the planned
out-of-sample evaluation is designed to fill.

%=========================================================
\appendix
\newpage
%=========================================================
\section{Foundations}\label{app:orig}
%=========================================================
This appendix derives, from the reduced-form VAR, the objects that the new results
of Appendix~\ref{app:new} take as given: the companion form and forward recursion,
the baseline forecast, the moving-average coefficients and their stacked form
$\Mmap$, the Waggoner--Zha minimum-norm solution $\hat u=\pinv{R}r$, and the
structural identification $u=D\eps$. Rather than reproduce the source paper's
results wholesale, we keep only what the extensions build on, and each subsection
closes with a pointer (\emph{Used in:}) to where it is used downstream. The source
paper's variance decomposition and out-of-sample metrics are recorded at the end
(\S\ref{app:refresults}) for reference; they are not prerequisites for the new
results. The structural-space restriction used in the main text, $R\beps=r$ with
stacked map $\Mmap$, is the reduced-form restriction here under $u=D\eps$
(\S\ref{app:norm}); the standalone \texttt{derivation\_companion.tex} carries the
same material in the source paper's own numbering and prose.

\medskip
%======================================================================
% Appendix A -- Foundations. Written to be \input by paper.tex (its \section
% header and intro live there). These are the derivations the NEW results of
% Appendix B take as given; each subsection ends with a pointer to where it is
% used. The source paper's variance decomposition and OOS metrics are recorded
% at the end for reference (they are not prerequisites for the new results).
%======================================================================

\subsection{Reduced form and companion form}\label{app:redcomp}
The starting point is the reduced-form VAR
\begin{equation}
y_t=\nu+A_1y_{t-1}+A_2y_{t-2}+\cdots+A_py_{t-p}+u_t,
\qquad \E[u_t]=0,\ \ \E[u_tu_t\tp]=\Sigma_u,
\tag{1}\label{aeq:red}
\end{equation}
a $K$-dimensional linear autoregression of order $p$. ``Reduced form'' means
$u_t$ collects statistical innovations with \emph{no} economic interpretation;
the $A_i$ summarise own- and cross-dynamics. Because every equation shares the
same right-hand-side regressors, the seemingly-unrelated-regression system
collapses to equation-by-equation OLS.

\paragraph{Companion (state-space) form.} Stacking the current and $p-1$ most
recent vectors turns \eqref{aeq:red} into a first-order system,
\begin{equation}
Y_t=\mu+AY_{t-1}+U_t,
\tag{2}\label{aeq:comp}
\end{equation}
with
\[
Y_t=\begin{bmatrix} y_t\\ y_{t-1}\\ \vdots\\ y_{t-p+1}\end{bmatrix},\quad
\mu=\begin{bmatrix}\nu\\0\\\vdots\\0\end{bmatrix},\quad
U_t=\begin{bmatrix}u_t\\0\\\vdots\\0\end{bmatrix},\quad
A=\begin{bmatrix}
A_1 & A_2 & \cdots & A_{p-1} & A_p\\
I_K & 0 & \cdots & 0 & 0\\
0 & I_K & \cdots & 0 & 0\\
\vdots & & \ddots & & \vdots\\
0 & 0 & \cdots & I_K & 0
\end{bmatrix}.
\]
Reading \eqref{aeq:comp} block-row by block-row: the first block row is exactly
\eqref{aeq:red}; block row $i\ge2$ is the identity $y_{t-i+1}=y_{t-i+1}$, which
merely carries lags forward in the state. With the top-block selector
$J=[\,I_K\ 0\ \cdots\ 0\,]\in\R^{K\times Kp}$ one has the identities used
repeatedly below,
\begin{equation}
y_t=JY_t,\qquad \mu=J\tp\nu,\qquad U_t=J\tp u_t.
\tag{$\star$}\label{aeq:star}
\end{equation}
Covariance-stationarity requires all eigenvalues of $A$ to lie inside the unit
circle.

\medskip
\noindent\textit{Used in:} everything below; the companion matrix $A$ is exported
to the dashboard so the counterfactual baseline (\S\ref{app:cf}) can be iterated
in the browser from any in-sample origin.

\subsection{Forward iteration: the \texorpdfstring{$h$}{h}-step recursion}\label{app:recursion}
Iterating \eqref{aeq:comp} forward $h$ periods gives
\begin{equation}
Y_{t+h}=\Bigg(\sum_{j=0}^{h-1}A^{j}\Bigg)\mu+A^{h}Y_t+\sum_{j=0}^{h-1}A^{j}U_{t+h-j}.
\tag{3}\label{aeq:recursion}
\end{equation}

\paragraph{Derivation (induction on $h$).} The base case $h=1$ is
\eqref{aeq:comp}. Assuming \eqref{aeq:recursion} at horizon $h$ and applying
\eqref{aeq:comp} once more, $Y_{t+h+1}=\mu+AY_{t+h}+U_{t+h+1}$, substitution gives
\[
Y_{t+h+1}
=\underbrace{\mu+\sum_{j=0}^{h-1}A^{j+1}\mu}_{=\sum_{j=0}^{h}A^{j}\mu}
+A^{h+1}Y_t
+\underbrace{\sum_{j=0}^{h-1}A^{j+1}U_{t+h-j}+U_{t+h+1}}_{=\sum_{j=0}^{h}A^{j}U_{t+h+1-j}},
\]
where the last bracket re-indexes $j\mapsto j+1$ (turning the sum into
$\sum_{j=1}^{h}A^{j}U_{t+h+1-j}$) and the stray $U_{t+h+1}$ supplies the missing
$j=0$ term. This is \eqref{aeq:recursion} at $h+1$.

\medskip
\noindent\textit{Used in:} the baseline (\S\ref{app:base}) is its conditional
mean; the moving-average map (\S\ref{app:ma}) is its stochastic part.

\subsection{The baseline forecast}\label{app:base}
Taking $\E[\,\cdot\mid\Fset_t]$ of \eqref{aeq:recursion} annihilates every future
shock, since $U_{t+h-j}$ is dated $t+h-j\ge t+1$ for $j=0,\dots,h-1$ and hence
unforecastable, $\E[U_{t+h-j}\mid\Fset_t]=0$. Premultiplying by $J$ and using
$\mu=J\tp\nu$ from \eqref{aeq:star},
\begin{equation}
\E[y_{t+h}\mid\Fset_t]
=\Bigg(\sum_{j=0}^{h-1}JA^{j}J\tp\Bigg)\nu+JA^{h}Y_t.
\tag{4}\label{aeq:base}
\end{equation}
Under quadratic loss this conditional expectation is the optimal point forecast;
it is the \emph{baseline scenario} against which every conditional forecast and
every counterfactual is measured.

\medskip
\noindent\textit{Used in:} the constraint right-hand side
$r=(\text{target}-\text{baseline})$ of the main text~\eqref{eq:restriction} and the counterfactual
origin baseline $F_{\text{base}}$ (\S\ref{app:cf}).

\subsection{Forecast error, impulse responses, and the stacked map}\label{app:ma}
Subtracting \eqref{aeq:base} from \eqref{aeq:recursion} leaves only the shock
term; projecting onto the top block with $U_{t+h-j}=J\tp u_{t+h-j}$ gives the
deviation of any realised path from the baseline,
\begin{equation}
y_{t+h}-\E[y_{t+h}\mid\Fset_t]
=\sum_{j=0}^{h-1}JA^{j}J\tp\,u_{t+h-j}
=\sum_{j=0}^{h-1}\Phi_j\,u_{t+h-j},
\qquad \Phi_j:=JA^{j}J\tp.
\tag{5}\label{aeq:dev}
\end{equation}
The $\Phi_j$ are the reduced-form moving-average (impulse-response) coefficients,
with $\Phi_0=I_K$. \textbf{Economic reading:} forcing a variable onto a scenario
path (left-hand side $\neq0$) is \emph{equivalent} to imposing a particular
realisation of future innovations (right-hand side); the scenario cannot happen
``for free,'' and the same innovations move every other variable through $\Phi_j$.

\paragraph{The stacked map.} Collect \eqref{aeq:dev} for $h=1,\dots,H$. Writing
$\Delta y=[\,(y_{t+1}-\E)\tp,\dots,(y_{t+H}-\E)\tp\,]\tp\in\R^{KH}$ and stacking
the innovations conformably, the deviation is a single linear image of the future
shocks,
\begin{equation}
\Delta y=\Mmap_u\,u,\qquad
\Mmap_u=\begin{bmatrix}
\Phi_0 & & &\\
\Phi_1 & \Phi_0 & &\\
\vdots & & \ddots &\\
\Phi_{H-1} & \Phi_{H-2} & \cdots & \Phi_0
\end{bmatrix},
\tag{5$'$}\label{aeq:stack}
\end{equation}
a block-lower-triangular $KH\times KH$ matrix (block $(h,s)=\Phi_{h-s}$ for
$s\le h$, zero otherwise). Once a structural impact matrix $D$ is identified
(\S\ref{app:ident2}), the same construction with $\Theta_j:=\Phi_jD$ in place of
$\Phi_j$ gives the \emph{structural} stacked map $\Mmap$ of the main text
\eqref{eq:Mmap}, whose diagonal block is $\Theta_0=D$.

\medskip
\noindent\textit{Used in:} $\Mmap$ is \emph{the} object of the paper. The
restriction rows $R$ are selected rows of $\Mmap$; the propagation
$\Delta y=\Mmap\beps$ delivers every conditional forecast, band
(\S\ref{app:shock}, \S\ref{app:param}), and counterfactual (\S\ref{app:cf}). The
coefficients $\Theta_j$ are exported to the dashboard, which rebuilds $\Mmap$ in
the browser.

\subsection{The conditioning restriction and the minimum-norm shock path}\label{app:minnorm}
A scenario fixes target values for a subset of variables over some horizons.
Reading the corresponding rows off \eqref{aeq:stack} yields a linear system in the
stacked future innovations,
\begin{equation}
R\,u=r,
\tag{5$''$}\label{aeq:Rur}
\end{equation}
where $r\in\R^{m}$ is the gap between the scenario values and the baseline
forecasts on the $m$ constrained cells and $R\in\R^{m\times KH}$ collects the
associated rows of the MA coefficients. Because a scenario fixes far fewer cells
than there are free shocks ($m<KH$), \eqref{aeq:Rur} is underdetermined; WZ, after
Doan--Litterman--Sims, select the minimum-norm solution.

\paragraph{Derivation (constrained minimisation).} Solve
$\min_u\tfrac12u\tp u$ s.t.\ $Ru=r$. The Lagrangian
$\mathcal L=\tfrac12u\tp u+\lambda\tp(r-Ru)$ has first-order conditions
$u=R\tp\lambda$ and $Ru=r$; substituting, $RR\tp\lambda=r$, so
$\lambda=(RR\tp)^{-1}r$ (full row rank $m$) and
\begin{equation}
\hat u=R\tp(RR\tp)^{-1}r=\pinv{R}\,r.
\tag{5$'''$}\label{aeq:minnorm}
\end{equation}
Adding $\hat u$ back through \eqref{aeq:stack} delivers the conditional forecast
for every variable.

\begin{remark}[dimensions of the inverse]
With $R$ of size $m\times KH$ and $m<KH$, $RR\tp$ ($m\times m$) is invertible
while $R\tp R$ ($KH\times KH$) is singular. A source-paper display that writes
$\hat u=R\tp(R\tp R)^{-1}r$ should read $(RR\tp)^{-1}$ as in \eqref{aeq:minnorm};
the two coincide only under the transposed storage convention for $R$. The
substance is unchanged.
\end{remark}

\begin{remark}[``most likely'' shocks]
If $u\sim N(0,I_{KH})$ then \eqref{aeq:minnorm} is the mode (equivalently the mean)
of $u\mid Ru=r$: minimising $u\tp u$ maximises the Gaussian density. This is why
$\hat u$ is the most likely combination of shocks consistent with the scenario ---
the reading the uncertainty derivations of Appendix~\ref{app:new} make precise.
\end{remark}

\medskip
\noindent\textit{Used in:} \eqref{aeq:minnorm} is the object every new derivation
extends --- soft conditioning replaces $(RR\tp)^{-1}$ by $(RR\tp+\Om)^{-1}$
(\S\ref{app:soft}); the shock band conditions the Gaussian around $\hat\beps$
(\S\ref{app:shock}); the counterfactual perturbs a realised $\hat u$
(\S\ref{app:cf}). The structural-space version $\hat\beps=\pinv{R}r$ used in the
main text is \eqref{aeq:minnorm} with $u=D\eps$ (\S\ref{app:norm}).

\subsection{Structural identification}\label{app:ident2}
The reduced-form innovation is decomposed as $u_t=D\eps_t$ with orthonormal
structural shocks,
\begin{equation}
\E[\eps_t]=0,\quad \E[\eps_t\eps_t\tp]=I_K
\ \Longrightarrow\ \Sigma_u=DD\tp.
\tag{6}\label{aeq:struct}
\end{equation}

\paragraph{Order condition and Cholesky.} A symmetric $\Sigma_u$ has $K(K+1)/2$
distinct entries while $D$ has $K^2$; the equations $\Sigma_u=DD\tp$ fall short by
$K^2-K(K+1)/2=K(K-1)/2$ restrictions. Requiring $D$ lower-triangular zeros out
exactly its $K(K-1)/2$ upper entries, closing the gap, and with a positive
diagonal the solution is the unique Cholesky factor
$D=\chol(\hat\Sigma_u)$, with $\hat\Sigma_u=\sum_t\hat u_t\hat u_t\tp/(T-Kp-1)$. A
lower-triangular $D$ means shock $k$ moves variables $1,\dots,k-1$ only with a
lag but $k,\dots,K$ contemporaneously, so the recursive ordering of $y_t$ is
material (global variables first, sticky prices ahead of the activity/policy
block, forward-looking financial variables last). The structural impulse
responses are $\Theta_j=\Phi_jD$.

\medskip
\noindent\textit{Used in:} $D=\Theta_0$ is the diagonal block of $\Mmap$ and the
change of variables $u=D\eps$ that turns the WZ norm into the Mahalanobis norm
(\S\ref{app:norm}); the whole apparatus is identification-agnostic --- replace
the Cholesky $D$ by a sign-restricted or instrument-based $D$ and every formula in
Appendix~\ref{app:new} is unchanged.

\subsection{Reference: variance decomposition and out-of-sample metrics}\label{app:refresults}
The following source-paper results are recorded for completeness; they are
\emph{not} prerequisites for the new derivations of Appendix~\ref{app:new}.

\paragraph{Forecast-error variance decomposition.} With the structural forecast
error $y_{t+h}-\E y_{t+h}=\sum_{j=0}^{h-1}\Theta_j\eps_{t+h-j}$ and orthonormal
$\eps$, the mean squared prediction error and its per-shock split are
\begin{equation}
\MSPE(h)=\sum_{j=0}^{h-1}\Theta_j\Theta_j\tp,\qquad
\MSPE_{k,l}(h)=\sum_{j=0}^{h-1}\theta_{k,l,j}^2,
\tag{7$'$,\,7$''$}\label{aeq:fevd}
\end{equation}
where $\theta_{k,l,j}=[\Theta_j]_{kl}$; the reported share is
$\MSPE_{k,l}(h)/\sum_{l}\MSPE_{k,l}(h)$, and orthogonality is what makes the shares
sum to one across $l$.

\paragraph{Out-of-sample evaluation (planned exercise).} Point accuracy uses an
AR($p$) benchmark --- \eqref{aeq:red} with diagonal $A_i$ --- and $h$-step errors
$e^{(h,m)}_{k,t}$ aggregated into
\begin{equation}
\RMSE(k,h,m)=\sqrt{\tfrac{1}{\bar t-\underline t}\textstyle\sum_t\big(e^{(h,m)}_{k,t}\big)^2},
\qquad
\MAE(k,h,m)=\tfrac{1}{\bar t-\underline t}\textstyle\sum_t\big|e^{(h,m)}_{k,t}\big|,
\tag{9,\,10}\label{aeq:rmse}
\end{equation}
each reported relative to the benchmark (a ratio below one favours the VAR).
Density accuracy uses the continuous ranked probability score in its quantile
(``pinball'') form: with $\rho_\tau(u)=u(\tau-\Id(u<0))$ and predictive quantiles
$\hat q^{(h,m)}_{\tau,k,t}$ on a grid $\tau_1<\cdots<\tau_N$,
\begin{equation}
\CRPS(k,h,m)=\frac{2}{N-1}\sum_{n=1}^{N}
\Bigg(\frac{1}{\bar t-\underline t}\sum_{t}\rho_{\tau_n}\!\big(y_{k,t}-\hat q^{(h,m)}_{\tau_n,k,t}\big)\Bigg),
\tag{11}\label{aeq:crps}
\end{equation}
the trapezoidal discretisation of
$\CRPS(F,y)=2\int_0^1\rho_\tau(y-F^{-1}(\tau))\,d\tau$, which needs only the
predictive quantiles, not the full density. These metrics drive the planned
backtesting harness (\S\ref{sec:oos}).

\subsection*{Summary of the foundation chain}
\[
\underbrace{(1)\Rightarrow(2)}_{\text{companion}}
\ \Rightarrow\
\underbrace{(3)}_{\text{iterate}}
\ \Rightarrow\
\begin{cases}
\E[\cdot\mid\Fset_t]\ \Rightarrow\ (4)\ \text{baseline}\\[2pt]
\text{deviation}\ \Rightarrow\ (5),(5')\ \Rightarrow\ \Mmap\ \text{(stacked map)}\\[2pt]
Ru=r\ \Rightarrow\ \hat u=\pinv{R}r\ \text{(min-norm, WZ)}\\[2pt]
u=D\eps,\ \Sigma_u=DD\tp\ \Rightarrow\ \Theta_j=\Phi_jD\ \text{(identification)}
\end{cases}
\]
Appendix~\ref{app:new} builds on these objects: it adds a distribution around
$\hat u$ (bands), a noisy version of $Ru=r$ (soft conditions), extra rows (joint /
partial-horizon scenarios), and a perturbation of a realised $\hat u$
(counterfactuals).

%=========================================================
\section{New derivations}\label{app:new}
%=========================================================

\subsection{Structural versus reduced-form norm}\label{app:norm}
Reduced-form innovations satisfy $u=D\eps$, stacked as
$\Delta y=\Mmap_u u=\Mmap\beps$ with $\Mmap_u$ the reduced-form stacked map (blocks
$\Phi_j$) and $\Mmap$ its structural counterpart (blocks $\Theta_j=\Phi_jD$). The
reduced-form restriction $R_u u=r$ and the structural one $R\beps=r$ have
$R=R_uD_{\text{blk}}$, $u=D_{\text{blk}}\beps$, where $D_{\text{blk}}=I_H\otimes D$.
Minimizing $\beps\tp\beps$ subject to $R\beps=r$ is minimizing
$u\tp D_{\text{blk}}^{-\top}D_{\text{blk}}^{-1}u=u\tp(I_H\otimes\Sigma_u^{-1})u$,
the Mahalanobis norm. Under $\beps\sim N(0,I)$ the density is
$\propto\exp(-\tfrac12\beps\tp\beps)$, so \eqref{eq:wzsol} is the mode of
$\beps\mid R\beps=r$, justifying ``most likely shocks.'' The two norms differ in
every unconstrained variable; the structural (Mahalanobis) norm is the correct one.

\subsection{The scenario-to-units map}\label{app:scenmap}
Let $\ell_0$ be the last observed level of the variable and $x_0$ its value. The
constructors used to author a path (all in \emph{level} space, then mapped by
\eqref{eq:units}) are:
\begin{itemize}[nosep,leftmargin=1.4em]
\item \emph{ramp} to $\bar\ell$ over $n$ months:
$\ell_h=x_0+(\bar\ell-x_0)\min(h/n,1)$;
\item \emph{hold} at $a$: $\ell_h=a$ for all $h$;
\item \emph{pulse} to $\bar\ell$ over $n$, hold $k$: the ramp-hold on
$h\le n{+}k$, and $\ell_h=\text{NaN}$ (free) for $h>n{+}k$;
\item \emph{policy}: $r_h$ steps up by $s$ each month to a peak $\bar r$, holds
$k$ months, then eases by $s$ to a floor $\underline r$ (a level-variable path);
\item \emph{replay} from date $\tau$ for $n$ months: copy the variable's own
realized growth (for a log-diff variable) or level (for a level variable) starting
just after $\tau$, then free.
\end{itemize}
A \emph{window} $[h_1,h_2]$ sets $r^{\text{unit}}_h=\text{NaN}$ outside the window;
NaN cells contribute no row to \eqref{eq:restriction}. For a log-difference
variable, note the asymmetry: constraining months $1,\dots,n$ pins the level at
month $n$, whereas a window that does not start at $h=1$ constrains growth rates
without pinning the level they build on.

\subsection{Joint and partial-horizon constraints}\label{app:joint}
For $M$ constrained variables with per-variable, per-horizon targets, the row set
of \eqref{eq:restriction} is
\[
\{\,(h,v): \text{target}(v,h)\ \text{is not NaN}\,\},\qquad
\text{row}(h,v)=(h-1)K+v,
\]
so $R$ is $\Mmap$ restricted to those rows and $r$ the corresponding
target-minus-baseline entries. Nothing else changes: the solve \eqref{eq:wzsol},
the soft update \eqref{eq:softpost}, and the bands \eqref{eq:shockpost} all take
this $R$ as input. Duplicated (variable, horizon) cells are disallowed (a variable
cannot be constrained twice at the same horizon).

\subsection{Soft conditioning (Gaussian update)}\label{app:soft}
Take the prior $\beps\sim N(0,I_{KH})$ and the noisy measurement
\eqref{eq:soft}, $r=R\beps+e$, $e\sim N(0,\Om)$, $e\perp\beps$. The joint law is
Gaussian with
\[
\Cov(\beps,r)=R\tp,\qquad \Var(r)=RR\tp+\Om,\qquad \E[\beps]=0,\ \E[r]=0.
\]
The standard Gaussian conditioning formula
$\E[\beps\mid r]=\Cov(\beps,r)\Var(r)^{-1}r$ and
$\Var[\beps\mid r]=\Var(\beps)-\Cov(\beps,r)\Var(r)^{-1}\Cov(r,\beps)$ give
\eqref{eq:softpost}:
\[
\hat\beps_{\text{soft}}=R\tp(RR\tp+\Om)^{-1}r,\qquad
P_{\text{soft}}=I-R\tp(RR\tp+\Om)^{-1}R.
\]
As $\Om\to0$, $R\tp(RR\tp)^{-1}=\pinv{R}$ and $P_{\text{soft}}\to I-\pinv{R}R=P$,
the hard-condition projector. A posterior draw is generated without forming
$P_{\text{soft}}$: with $z\sim N(0,I)$ and $e\sim N(0,\Om)$,
\[
\beps^{\text{draw}}=\hat\beps_{\text{soft}}+\big(z-R\tp(RR\tp+\Om)^{-1}(Rz+e)\big),
\]
whose mean is $\hat\beps_{\text{soft}}$ and whose covariance is exactly
$P_{\text{soft}}$; $\Om=0$ collapses it to the hard projector draw of
\S\ref{app:shock}.

\subsection{Shock-uncertainty band}\label{app:shock}
Condition $\beps\sim N(0,I)$ on $R\beps=r$ (hard). Decompose $\beps$ into its
component in $\mathrm{row}(R)$ and its component in $\ker R$. The row-space
component is fixed by the constraint at $\pinv{R}r$; the kernel component is
unconstrained and retains its prior law. Hence
$\beps\mid R\beps=r\sim N(\hat\beps,P)$ with $\hat\beps=\pinv{R}r$ and
$P=I-\pinv{R}R$ the orthogonal projector onto $\ker R$ (idempotent, symmetric,
$RP=0$). A draw is $\hat\beps+Pz=\hat\beps+(z-\pinv{R}(Rz))$, applied through the
shared factors of $R$ so $P$ is never formed. Propagating,
$\Delta y=\Mmap\beps\sim N(\Mmap\hat\beps,\Mmap P\Mmap\tp)$, which is
\eqref{eq:shockband}. Since $RP=0$, the constrained rows of $\Mmap P\Mmap\tp$
vanish and the constrained cells have zero predictive variance---the pinch in the
fan chart. Because the display transform (level reconstruction $\ell_0\exp(\sum)$,
year-over-year growth) is nonlinear, bands are read as sample quantiles of the
transformed draws, not as transforms of the Gaussian quantiles.

\subsection{Parameter-uncertainty band}\label{app:param}
Let $\{(A^{(d)},\Sigma_u^{(d)})\}_{d=1}^{N}$ be draws of the VAR parameters.
For each: form $D^{(d)}=\chol(\Sigma_u^{(d)})$, the structural map $\Mmap^{(d)}$
(blocks $\Theta_j^{(d)}=\Phi_j^{(d)}D^{(d)}$), the baseline $F^{(d)}_{\text{base}}$,
then rebuild $R^{(d)},r^{(d)}$ and solve $\hat\beps^{(d)}=(R^{(d)})^{+}r^{(d)}$.
The per-draw conditional \emph{mean} is
$F^{(d)}_{\text{base}}+\Mmap^{(d)}\hat\beps^{(d)}$ and a full predictive draw adds
a shock draw $\Mmap^{(d)}(z-\text{proj})$ from \S\ref{app:shock} at draw $d$'s
parameters. Reported quantities are quantiles, across $d$, of the transformed
conditional draws (\emph{cond}); of the baseline draws (\emph{base}); and of the
scenario-minus-baseline means $\Mmap^{(d)}\hat\beps^{(d)}$ (\emph{diff}), which
isolate parameter uncertainty because with parameters fixed the difference is
deterministic.

\paragraph{Draw mechanism.} With a Minnesota prior, draws are from the
normal--inverse--Wishart posterior (rejecting explosive companions); otherwise a
residual bootstrap regenerates a sample and re-estimates. For the bootstrap, the
data-generating companion is bias-adjusted following Kilian (1998): estimate the
mean bias by a first bootstrap pass, subtract it, and shrink the adjustment toward
zero until the companion is nonexplosive; residuals are resampled only from months
not activated by an impulse dummy (whose residuals are zero by construction), so
those months are not smuggled back in. Explosive draws in the main pass are
redrawn. The OLS baseline is \emph{not} the bootstrap median (the bias adjustment
and rejection shift the per-draw baseline), which is why the \emph{diff} object,
computed draw by draw, is the correct basis for ``distinguishable from baseline.''

\subsection{Historical counterfactual}\label{app:cf}
Fix an origin $t_0$ within the sample and a window of length $H$. The realized
data over the window equal the origin baseline plus the realized structural
innovations, $y_{\text{actual}}=F_{\text{base}}+\Mmap\beps_{\text{actual}}$, so on
the constrained rows $R\beps_{\text{actual}}=r_{\text{actual}}$ with
$r_{\text{actual}}=(y_{\text{actual}}-F_{\text{base}})_{[\text{rows}]}$. Impose a
counterfactual target $r_{\text{cf}}$ on those rows and keep every other realized
shock by minimizing the change:
\[
\min_{\delta}\lVert\delta\rVert\ \text{ s.t. }\ R(\beps_{\text{actual}}+\delta)=r_{\text{cf}}
\ \Rightarrow\ R\delta=r_{\text{cf}}-r_{\text{actual}}
\ \Rightarrow\ \delta=\pinv{R}(r_{\text{cf}}-r_{\text{actual}}),
\]
by the same Lagrangian argument as \eqref{eq:wzsol}. Then
$y_{\text{cf}}=y_{\text{actual}}+\Mmap\delta$, i.e.\ \eqref{eq:cf2}; the baseline
cancels because it is common to actual and counterfactual. To confine the change
to a subset $S$ of structural shocks, restrict $\delta$ to the columns of $R$
indexed by $S$ across horizons, $R_S=R_{:,[S]}$, solve
$\delta_S=\pinv{(R_S)}(r_{\text{cf}}-r_{\text{actual}})$, and embed $\delta_S$ back
into a zero vector. The reported perturbation size is
$\lVert\delta\rVert/\sqrt{KH}$; a rank-deficient $R_S$ (an infeasible confined
target) is handled by the pseudoinverse, which projects the shortfall away. The
reduced-form VAR is held fixed throughout---the Lucas-critique assumption that
bounds the interpretation.


%=========================================================
\begin{thebibliography}{9}
\bibitem{bm2017} Blake, A.\ and Mumtaz, H.\ (2017). \emph{Applied Bayesian
Econometrics for Central Bankers}. Bank of England, CCBS Technical Handbook.
\bibitem{dls1984} Doan, T., Litterman, R.\ and Sims, C.\ (1984). Forecasting and
conditional projection using realistic prior distributions. \emph{Econometric
Reviews} 3, 1--100.
\bibitem{kilian1998} Kilian, L.\ (1998). Small-sample confidence intervals for
impulse response functions. \emph{Review of Economics and Statistics} 80, 218--230.
\bibitem{mss2025} Moran, K., Stevanovic, D.\ and Surprenant, S.\ (2025). Risk
Scenarios and Macroeconomic Forecasts. Bank of Canada Staff Working Paper 2025-28.
\bibitem{wz1999} Waggoner, D.\ and Zha, T.\ (1999). Conditional forecasts in
dynamic multivariate models. \emph{Review of Economics and Statistics} 81, 639--651.
\bibitem{adrr2018} Antol\'in-D\'iaz, J.\ and Rubio-Ram\'irez, J.\ (2018). Narrative sign restrictions for SVARs. \emph{American Economic Review} 108, 2802--2829.
\bibitem{gk2015} Gertler, M.\ and Karadi, P.\ (2015). Monetary policy surprises, credit costs, and economic activity. \emph{American Economic Journal: Macroeconomics} 7, 44--76.
\bibitem{mr2013} Mertens, K.\ and Ravn, M.\ (2013). The dynamic effects of personal and corporate income tax changes in the United States. \emph{American Economic Review} 103, 1212--1247.
\bibitem{rigobon2003} Rigobon, R.\ (2003). Identification through heteroskedasticity. \emph{Review of Economics and Statistics} 85, 777--792.
\bibitem{rwz2010} Rubio-Ram\'irez, J., Waggoner, D.\ and Zha, T.\ (2010). Structural vector autoregressions: theory of identification and algorithms for inference. \emph{Review of Economic Studies} 77, 665--696.
\bibitem{sw2018} Stock, J.\ and Watson, M.\ (2018). Identification and estimation of dynamic causal effects in macroeconomics using external instruments. \emph{Economic Journal} 128, 917--948.
\bibitem{uhlig2005} Uhlig, H.\ (2005). What are the effects of monetary policy on output? Results from an agnostic identification procedure. \emph{Journal of Monetary Economics} 52, 381--419.
\end{thebibliography}

\end{document}
